% GENERATED FILE. Edit canonical lessons, then run npm run generate:books.
% canonical-source-hash: fnv1a64:6ef004024a0af93c
% canonical-lessons: AR-C459-mustahlik, AR-C459-mizaniyya, AR-C459-masarif, AR-C459-mushtarayat, AR-C459-iqtasada

\chapter{Spending wisely}
\label{ch:ar-spending-wisely}
\hlchaptermodality{459}

\begin{chapteropening}
\textbf{By the end of this chapter:} \emph{I can say a consumer, a budget, expenses, purchases and to economise.}

Complete the last lesson of chapter 459: I can say a consumer, a budget, expenses, purchases and to economise.
\end{chapteropening}

\section[\texorpdfstring{\ar{مستهلك} (mustahlik)}{مستهلك (mustahlik)}]{\ar{مستهلك} (mustahlik) — a consumer}

\label{lesson:AR-C459-mustahlik}



\begin{quote}
Before the new one: say the Arabic for to take back, then the Arabic for to barter.
\end{quote}

\subsection*{You'll want to know: \ar{مستهلك}}
\textbf{\ar{مستهلك}} — \emph{mustahlik} — \textquotedblleft{}a consumer\textquotedblright{}.

\begin{quote}
\ar{يحب} \ar{المستهلك} \ar{الأسعار} \ar{الرخيصة}.
\end{quote}

\begin{quote}
\emph{yuḥibbu al-mustahlik al-asʿār ar-rakhīṣa.}
\end{quote}

\begin{quote}
The consumer likes cheap prices.
\end{quote}

\begin{cousinweb}
\textbf{\ar{مستهلك}} (\emph{mustahlik}) is built on the native Arabic root \emph{h-l-k}, \textquotedblleft{}to perish\textquotedblright{} (Lane). It is the Form X participle, \textquotedblleft{}one who uses things up\textquotedblright{}.
\end{cousinweb}

\begin{grammarlens}[title={What you've built}]
One more word for getting around.
\end{grammarlens}

\subsection*{Your turn}
\begin{itemize}
\raggedright
  \item \emph{Say it:} \emph{mustahlik}
  \item \emph{Say it:} \emph{mustahlik}, once more
  \item \emph{Say it:} \emph{yuḥibbu al-mustahlik al-asʿār ar-rakhīṣa.} — The consumer likes cheap prices.
  \item \emph{Recall:} say the Arabic for to take back, then the Arabic for to barter, then say \emph{mustahlik} again
\end{itemize}

\begin{tcolorbox}[breakable,colback=teal!4,colframe=teal!35!black,title={Before you move on}]
What does \ar{مستهلك} mean? (\textquotedblleft{}A consumer\textquotedblright{}.) Say it once more.
\end{tcolorbox}
\section[\texorpdfstring{\ar{ميزانية} (mīzāniyya)}{ميزانية (mīzāniyya)}]{\ar{ميزانية} (mīzāniyya) — a budget}

\label{lesson:AR-C459-mizaniyya}



\begin{quote}
Before the new one: say the Arabic for to barter, then the Arabic for a consumer.
\end{quote}

\subsection*{You'll want to know: \ar{ميزانية}}
\textbf{\ar{ميزانية}} — \emph{mīzāniyya} — \textquotedblleft{}a budget\textquotedblright{}.

\begin{quote}
\ar{كانت} \ar{ميزانية} \ar{الرحلة} \ar{صغيرة}.
\end{quote}

\begin{quote}
\emph{kānat mīzāniyyat ar-riḥla ṣaghīra.}
\end{quote}

\begin{quote}
The budget for the trip was small.
\end{quote}

\begin{cousinweb}
\textbf{\ar{ميزانية}} (\emph{mīzāniyya}) is built on the native Arabic root \emph{w-z-n}, \textquotedblleft{}to weigh\textquotedblright{} (Lane). It is made from \emph{mīzān}, \textquotedblleft{}scales\textquotedblright{}, with \emph{-iyya}: money weighed out.
\end{cousinweb}

\begin{grammarlens}[title={What you've built}]
One more word for getting around.
\end{grammarlens}

\subsection*{Your turn}
\begin{itemize}
\raggedright
  \item \emph{Say it:} \emph{mīzāniyya}
  \item \emph{Say it:} \emph{mīzāniyya}, once more
  \item \emph{Say it:} \emph{kānat mīzāniyyat ar-riḥla ṣaghīra.} — The budget for the trip was small.
  \item \emph{Recall:} say the Arabic for to barter, then the Arabic for a consumer, then say \emph{mīzāniyya} again
\end{itemize}

\begin{tcolorbox}[breakable,colback=teal!4,colframe=teal!35!black,title={Before you move on}]
What does \ar{ميزانية} mean? (\textquotedblleft{}A budget\textquotedblright{}.) Say it once more.
\end{tcolorbox}
\section[\texorpdfstring{\ar{مصاريف} (maṣārīf)}{مصاريف (maṣārīf)}]{\ar{مصاريف} (maṣārīf) — expenses}

\label{lesson:AR-C459-masarif}



\begin{quote}
Before the new one: say the Arabic for a consumer, then the Arabic for a budget.
\end{quote}

\subsection*{You'll want to know: \ar{مصاريف}}
\textbf{\ar{مصاريف}} — \emph{maṣārīf} — \textquotedblleft{}expenses\textquotedblright{}.

\begin{quote}
\ar{كانت} \ar{مصاريف} \ar{الفندق} \ar{كبيرة}.
\end{quote}

\begin{quote}
\emph{kānat maṣārīf al-funduq kabīra.}
\end{quote}

\begin{quote}
The hotel expenses were large.
\end{quote}

\begin{cousinweb}
\textbf{\ar{مصاريف}} (\emph{maṣārīf}) is built on the native Arabic root \emph{ṣ-r-f}, \textquotedblleft{}to turn away, to change\textquotedblright{} (Lane). Expenses are money turned out of the purse.
\end{cousinweb}

\begin{grammarlens}[title={What you've built}]
One more word for getting around.
\end{grammarlens}

\subsection*{Your turn}
\begin{itemize}
\raggedright
  \item \emph{Say it:} \emph{maṣārīf}
  \item \emph{Say it:} \emph{maṣārīf}, once more
  \item \emph{Say it:} \emph{kānat maṣārīf al-funduq kabīra.} — The hotel expenses were large.
  \item \emph{Recall:} say the Arabic for a consumer, then the Arabic for a budget, then say \emph{maṣārīf} again
\end{itemize}

\begin{tcolorbox}[breakable,colback=teal!4,colframe=teal!35!black,title={Before you move on}]
What does \ar{مصاريف} mean? (\textquotedblleft{}Expenses\textquotedblright{}.) Say it once more.
\end{tcolorbox}
\section[\texorpdfstring{\ar{مشتريات} (mushtarayāt)}{مشتريات (mushtarayāt)}]{\ar{مشتريات} (mushtarayāt) — purchases; the shopping}

\label{lesson:AR-C459-mushtarayat}



\begin{quote}
Before the new one: say the Arabic for a budget, then the Arabic for expenses.
\end{quote}

\subsection*{You'll want to know: \ar{مشتريات}}
\textbf{\ar{مشتريات}} — \emph{mushtarayāt} — \textquotedblleft{}purchases; the shopping\textquotedblright{}.

\begin{quote}
\ar{وضعت} \ar{المشتريات} \ar{في} \ar{السيارة}.
\end{quote}

\begin{quote}
\emph{waḍaʿtu al-mushtarayāt fī as-sayyāra.}
\end{quote}

\begin{quote}
I put the shopping in the car.
\end{quote}

\begin{cousinweb}
\textbf{\ar{مشتريات}} (\emph{mushtarayāt}) is built on the native Arabic root \emph{sh-r-y}, \textquotedblleft{}to buy, to sell\textquotedblright{} (Lane). It is a plural of the Form VIII passive participle: things bought.
\end{cousinweb}

\begin{grammarlens}[title={What you've built}]
One more word for getting around.
\end{grammarlens}

\subsection*{Your turn}
\begin{itemize}
\raggedright
  \item \emph{Say it:} \emph{mushtarayāt}
  \item \emph{Say it:} \emph{mushtarayāt}, once more
  \item \emph{Say it:} \emph{waḍaʿtu al-mushtarayāt fī as-sayyāra.} — I put the shopping in the car.
  \item \emph{Recall:} say the Arabic for a budget, then the Arabic for expenses, then say \emph{mushtarayāt} again
\end{itemize}

\begin{tcolorbox}[breakable,colback=teal!4,colframe=teal!35!black,title={Before you move on}]
What does \ar{مشتريات} mean? (\textquotedblleft{}Purchases; the shopping\textquotedblright{}.) Say it once more.
\end{tcolorbox}
\section[\texorpdfstring{\ar{اقتصد} (iqtaṣada)}{اقتصد (iqtaṣada)}]{\ar{اقتصد} (iqtaṣada) — to economise; to save}

\label{lesson:AR-C459-iqtasada}



\begin{quote}
Before the new one: say the Arabic for expenses, then the Arabic for purchases.
\end{quote}

\subsection*{You'll want to know: \ar{اقتصد}}
\textbf{\ar{اقتصد}} — \emph{iqtaṣada} — \textquotedblleft{}to economise; to save\textquotedblright{}.

\begin{quote}
\ar{اقتصدنا} \ar{في} \ar{المصاريف} \ar{هذا} \ar{الشهر}.
\end{quote}

\begin{quote}
\emph{iqtaṣadnā fī al-maṣārīf hādhā ash-shahr.}
\end{quote}

\begin{quote}
We economised on expenses this month.
\end{quote}

Present: \ar{يقتصد}, \emph{yaqtaṣid}.

\begin{cousinweb}
\textbf{\ar{اقتصد}} (\emph{iqtaṣada}) is built on the native Arabic root \emph{q-ṣ-d}, \textquotedblleft{}to aim at; to keep to the middle\textquotedblright{} (Lane). It is Form VIII, \textquotedblleft{}to be moderate\textquotedblright{}.
\end{cousinweb}

\begin{grammarlens}[title={What you've built}]
One more word for getting around.
\end{grammarlens}

\subsection*{Your turn}
\begin{itemize}
\raggedright
  \item \emph{Say it:} \emph{iqtaṣada}
  \item \emph{Say it:} \emph{iqtaṣada}, once more
  \item \emph{Say it:} \emph{iqtaṣadnā fī al-maṣārīf hādhā ash-shahr.} — We economised on expenses this month.
  \item \emph{Recall:} say the Arabic for expenses, then the Arabic for purchases, then say \emph{iqtaṣada} again
\end{itemize}

\begin{tcolorbox}[breakable,colback=teal!4,colframe=teal!35!black,title={Before you move on}]
What does \ar{اقتصد} mean? (\textquotedblleft{}To economise; to save\textquotedblright{}.) Say it once more.
\end{tcolorbox}
